% !PRINTED none
\begin{center}
\gesperrt{Fourth Chapter.}
\end{center}

\begin{center}
\gesperrt{Cubic and biquadratic Abelian fields.}
\end{center}

\begin{center}
\rule{3cm}{0.4pt}
\end{center}

\begin{center}
\S.\,23.
\end{center}

\begin{center}
\gesperrt{Cubic cyclotomic fields.}
\end{center}

The general investigations presented in the last sections admit manifold applications to special cases, of which the simplest shall be considered here. It is noteworthy in this connection that even the simplest cases, if one attacks them directly, lead almost into the same difficulties which we have overcome by the general investigation.

We shall first treat the problem of finding all \gesperrt{cubic cyclotomic fields, that is, all cyclotomic periods which satisfy a rational equation of degree $3$,} and indeed in such a way that, of the several expressions of which the same quantity is capable, the simplest is set down.

We have here only one invariant $i_1=3$, and according to the prescriptions in the summary of \S.\,21 there may be taken up into $m$ the factor 9, but not 3, and further primes in arbitrary number of the form $6N+1$, that is the primes 7, 13, 19, 31, 37, 43, 61, 67, 73, 79, 97 . . ., and so we thus obtain as suitable values of $m$ under 200:
\[
\begin{aligned}
m &= 7,\ 9,\ 13,\ 19,\ 31,\ 37,\ 43,\ 61,\ \mathbf{63},\ 67,\ 73,\ 79,\ \mathbf{91},\ 97,\ 103,\ 109,\\
&\quad \mathbf{117},\ 127,\ \mathbf{133},\ 139,\ 151,\ 157,\ 163,\ \mathbf{171},\ 181,\ 193,\ 199,
\end{aligned}
\]
among which the numbers with several prime divisors are distinguished by bold type. The smallest number $m$ which more than two primes divide is $7.9.13=819$.
